
\section{Evolutionary Code and Algorithm Discovery}
\label{sec:evolutionary-code-algorithm-discovery}

The recent literature on AI self-evolution contains a striking convergence: the object being improved is no longer only a model checkpoint, a prompt, or a policy, but an executable artifact that can be inspected, modified, evaluated, archived, and selected.  In this chapter, \emph{evolutionary code and algorithm discovery} denotes this family of methods.  The term covers prompt optimization systems such as OPRO, program-search systems such as FunSearch, codebase-level evolutionary agents such as AlphaEvolve and OpenEvolve, and agent-level systems such as ADAS, the Darwin G{\"o}del Machine, G{\"o}del Agent, Agent Symbolic Learning, Absolute Zero, and SelfEvolve.  These systems differ in scale and substrate, yet they instantiate the same closed loop: a population of candidate artifacts is represented in language or code; a generator proposes variants; an evaluator assigns scores; a memory or archive stores outcomes; and a selection rule decides what information should condition the next proposal.

This loop is central to self-evolution because it gives large language models (LLMs) a structured way to improve systems without changing the base model weights.  Standard LLM inference is episodic: the model receives a context, emits an answer, and forgets the attempt unless a user or external memory preserves it.  Evolutionary code discovery makes the episode persistent.  A failed program becomes a negative example, a successful heuristic becomes a parent, an error trace becomes mutation guidance, and a benchmark score becomes a scalar reward.  Over many iterations, the surrounding system can accumulate knowledge even when the LLM itself remains frozen.  This is why code evolution occupies a bridge position between prompt engineering, AutoML, reinforcement learning with verifiable rewards, and open-ended evolutionary computation.

The chapter follows six questions.  First, how did OPRO show that an LLM can operate as an optimizer through natural language history rather than through gradients?  Second, how can an LLM be treated as an evolutionary operator that performs mutation, crossover, and selection-aware recombination over programs?  Third, why did FunSearch make program evolution credible as a route to mathematical discovery, and how should its achievements be distinguished from the later AlphaEvolve result on matrix multiplication?  Fourth, what do open-source systems expose about engineering requirements, licenses, and reproducibility?  Fifth, how do agent self-evolution systems extend the same loop from individual functions to whole agent architectures?  Sixth, what benchmark evidence is currently available, and what are its methodological limitations?

A useful abstraction is to separate the \emph{inner solver} from the \emph{outer evolutionary process}.  The inner solver may be a prompt, a Python function, a scheduling heuristic, a reasoning trajectory, or a multi-agent workflow.  The outer process owns candidate generation, evaluation, archive maintenance, and search policy.  In classical evolutionary computation, mutation and crossover are usually syntactic or numeric operators.  In LLM-driven evolution, these operators become semantic: a model can read a parent program, infer its intent, identify a bottleneck, import a new algorithmic idea, and write a child that is not a small edit in token space but may still be a meaningful descendant in design space.  This semantic variation is the key advantage and the key danger.  It allows discontinuous jumps to powerful ideas, but it also makes safety, credit assignment, and reproducibility harder.

The strongest systems therefore combine neural creativity with hard verification.  OPRO evaluates candidate instructions on held-out tasks.  FunSearch executes programs in a sandbox and admits only candidates that pass syntactic and semantic filters.  AlphaEvolve and its open-source descendants rely on automated evaluators for objective scores.  SE-Agent evaluates evolved reasoning trajectories on SWE-bench Verified.  Absolute Zero uses a code executor to verify self-generated tasks and solutions.  The common lesson is that self-evolution is not merely ``asking the model to improve itself.''  It is the design of an evaluation ecology in which generated artifacts can safely compete, fail, and leave traces that shape future generations.

\subsection{LLM as Optimizer: OPRO}
\label{subsec:opro}

OPRO, or Optimization by PROmpting, introduced a clean formulation of the LLM-as-optimizer paradigm \cite{yang2023large}.  The central idea is deceptively simple.  Instead of hand-designing a gradient estimator, a Bayesian surrogate, or a genetic operator, the system expresses the optimization problem in natural language.  It gives the optimizer LLM a history of candidate solutions and their scores, asks for new candidates that should perform better, evaluates those candidates with a scorer, and repeats the process.  When the candidate is an instruction for a downstream LLM, OPRO becomes prompt optimization.  When the candidate is a vector of numerical variables or a route for a combinatorial problem, the same outer loop becomes a generic black-box optimizer.

Let $\mathcal{X}$ be a candidate space and let $f:\mathcal{X}\rightarrow\mathbb{R}$ be a black-box objective.  At iteration $t$, OPRO maintains a history
\begin{equation}
  \mathcal{H}_t = \{(x_i, f(x_i))\}_{i=1}^{n_t},
\end{equation}
where $x_i$ may be a natural-language instruction, a parameter tuple, or a combinatorial object serialized as text.  A meta-prompt constructor $M$ converts the history, task description $d$, constraints $c$, and optional examples $e$ into a prompt:
\begin{equation}
  m_t = M(\mathcal{H}_t, d, c, e).
\end{equation}
The optimizer LLM with parameters $\theta$ samples one or more new candidates,
\begin{equation}
  x_{t+1}^{(j)} \sim p_\theta(x \mid m_t, \tau_{\mathrm{opt}}), \qquad j=1,\ldots,k,
\end{equation}
where $\tau_{\mathrm{opt}}$ is typically high enough to maintain diversity.  After filtering duplicates or invalid outputs, the candidates are scored,
\begin{equation}
  s_{t+1}^{(j)} = f\bigl(x_{t+1}^{(j)}\bigr),
\end{equation}
and the history is updated by appending the new pairs and optionally truncating to the most informative top-scoring examples:
\begin{equation}
  \mathcal{H}_{t+1} = \operatorname{Truncate}\left(\mathcal{H}_{t}\cup\{(x_{t+1}^{(j)}, s_{t+1}^{(j)})\}_{j=1}^{k}\right).
\end{equation}
The final output is usually $x^*=\arg\max_{(x,s)\in\mathcal{H}_T}s$ for maximization, or $\arg\min$ for loss minimization.

This formulation is important because the optimizer is not trained on gradients of $f$.  It is given a textual trace of empirical evidence and asked to infer a better move.  The prompt therefore becomes an optimizer interface.  If the history is ordered, bucketed, or annotated, the model can exploit ordinal relationships.  If the prompt includes high-frequency mistakes, the model can target weak regions.  If the prompt includes constraints, the model can internalize feasibility.  The optimization algorithm is partly encoded in the surrounding program and partly encoded in the model's pretraining distribution over explanations, heuristics, and problem-solving patterns.

In prompt optimization, OPRO uses a two-layer architecture.  The outer layer is the meta-prompt supplied to the optimizer LLM.  It contains prior instructions and their scores, a statement that higher scores are better, and a request to generate a new instruction in a parseable format.  The inner layer is the task prompt supplied to the scorer LLM.  For each candidate instruction, the scorer LLM solves a collection of training examples, and the system computes an accuracy or task-specific score.  The local materials report four instruction placement modes: before the question, at the beginning of the question, at the end of the question, and at the beginning of the answer.  This matters because prompt optimization is not only the search for words; it is also the search for where those words interact with the model's generation policy.

The implementation details illustrate a general engineering pattern for self-evolution.  The optimizer temperature is high, encouraging diverse candidate generation, while the scorer temperature is low, making evaluation more deterministic.  OPRO's default instruction-optimization loop can run for about two hundred search steps and generate multiple candidate instructions per step.  The meta-prompt may retain a bounded number of top instructions rather than the entire history, both to fit within context and to emphasize successful patterns.  Scores can be discretized into buckets so that the optimizer does not overfit tiny score differences that may be noise.  Candidate extraction uses explicit delimiters such as \texttt{<INS>} and \texttt{</INS>} for one model family and bracket conventions for another; invalid, duplicate, overly long, or obviously malformed instructions are filtered before scoring.

The mathematical significance of OPRO is that it turns optimization into conditional generation over a history of evaluations.  The search policy is not a hand-coded evolutionary rule, yet it can emulate evolutionary behavior.  A high-scoring instruction acts as an elite.  A low-scoring instruction acts as a negative example.  A newly sampled instruction may mutate phrasing, recombine ideas from several parents, or introduce a new heuristic learned from pretraining.  Selection is explicit in the history truncation and ranking.  Thus OPRO is a minimal bridge from prompt engineering to evolutionary search: it lacks a population database, island model, or formal quality-diversity archive, but it already contains the essential ingredients of generate--evaluate--remember.

For reasoning tasks, OPRO's optimizer can rediscover familiar prompt patterns such as step-by-step reasoning, careful checking, or decomposition.  The point is not that every discovered instruction is semantically surprising.  The point is that the system can discover, rank, and reuse such instructions automatically for a specific task distribution.  This task adaptation is especially relevant in AI self-evolution because a deployed agent often faces a distribution that differs from the model's original training mixture.  An OPRO-like outer loop can tune the agent's operating prompt to the environment without requiring gradient updates.

OPRO also demonstrates that the same language interface can optimize non-linguistic variables.  In the linear-regression example, a candidate may be represented as a textual pair such as $w=15, b=14$, and the objective is a squared error,
\begin{equation}
  f(w,b) = \lVert y - (Xw+b)\rVert_2^2.
\end{equation}
In a traveling-salesperson example, a candidate is a route serialized as a trace and scored by tour length.  The optimizer LLM does not receive derivatives with respect to $w$, $b$, or the route.  It receives a record of previous values and objective scores.  This suggests a broad design principle: when a domain can be serialized into a representation that an LLM can interpret and when candidate quality can be automatically evaluated, an LLM can serve as a black-box proposal distribution over improvements.

However, OPRO also exposes limitations that persist in later systems.  First, the optimizer can be misled by noisy evaluations.  If the scorer LLM is stochastic or the training set is small, the history may reward spurious instructions.  Second, the representation of history is a bottleneck.  The model can only condition on what fits into context, so summarization and selection policies become part of the optimizer.  Third, OPRO's natural-language mutations have no formal guarantee of monotonic improvement.  Even when the prompt says ``generate a better instruction,'' the sample may be worse, redundant, or invalid.  Fourth, prompt optimization can overfit to the training examples used for scoring.  Validation intervals and held-out sets are therefore essential.

For a survey of self-evolution, the main contribution of OPRO is conceptual.  It shows that a language model can be wrapped as a general-purpose optimizer if the search state is expressed in language and if a reliable score is available.  Later systems replace instructions with functions, agents, trajectories, and codebases, but they preserve the same core equation: new artifacts are sampled from a model conditioned on the archive of old artifacts and their measured consequences.  OPRO is thus the first step in a progression from ``LLM writes an answer'' to ``LLM optimizes the process that writes answers.''

\subsection{LLM as Evolutionary Operator: Mutation, Crossover, and Selection}
\label{subsec:llm-evolutionary-operator}

Once an LLM is used to generate candidates from a history of scored artifacts, it is natural to interpret it as an evolutionary operator.  Classical evolutionary algorithms usually define mutation as a random perturbation, crossover as a recombination of parent genomes, and selection as a rule that biases reproduction toward higher-fitness individuals.  LLM-driven evolution changes the representation and the operators.  The genome may be a program, a prompt, a workflow, or a reasoning trajectory.  Mutation may be a semantic edit proposed after reading an error trace.  Crossover may be a synthesis of two algorithmic ideas described in code comments.  Selection may be implemented by an archive that exposes successful parents to the next LLM call while hiding or down-weighting failures.

A general LLM-evolution loop can be written as follows.  Let $A_t$ be an archive of artifacts $a$ with fitness $F(a)$ and optional behavior descriptor $b(a)$.  Parent selection chooses one or more parents,
\begin{equation}
  (a_{p_1},\ldots,a_{p_m})\sim q_t(a\mid A_t),
\end{equation}
where $q_t$ may depend on fitness, novelty, age, or niche coverage.  A prompt constructor serializes parent code, evaluation traces, failures, and instructions into $z_t$.  The LLM samples a child,
\begin{equation}
  a_{c}\sim p_\theta(a\mid z_t, \tau, \mathrm{mode}),
\end{equation}
where \textit{mode} may request mutation, recombination, repair, simplification, or targeted optimization.  The evaluator executes or judges $a_c$, producing $F(a_c)$ and $b(a_c)$.  The archive update rule then decides whether to insert, discard, or quarantine the child:
\begin{equation}
  A_{t+1}=U(A_t,a_c,F(a_c),b(a_c)).
\end{equation}
This formalism covers OPRO, FunSearch, AlphaEvolve, OpenEvolve, CodeEvolve, and agent-level variants.  The differences lie in $q_t$, the prompt constructor, the evaluator, and $U$.

LLM mutation is more powerful than token-level mutation because it can preserve intent while changing implementation.  For example, a model may replace a quadratic loop with a heap, introduce memoization, rewrite a heuristic priority function, or add a guard that prevents a known failure mode.  These changes may be large in edit distance but small in conceptual distance.  Conversely, a tiny textual edit may be semantically catastrophic.  This mismatch means that evolutionary distance in LLM systems should not be measured only by tokens or lines.  It should be measured by behavioral descriptors, test signatures, resource profiles, and genealogical relationships.

LLM crossover is similarly semantic.  A classical crossover operator splices substrings or subtrees.  A model-based operator can read two parents and produce a coherent child that combines the first parent's data structure with the second parent's scoring rule.  It can also reject a naive combination if the pieces are incompatible.  Systems such as SE-Agent make this explicit at the trajectory level: recombination takes useful fragments from multiple reasoning paths and constructs a new path intended to inherit their strengths.  Open-source code evolution systems present multiple ``inspiration'' programs to the model, which gives the model context for recombination without forcing a mechanical splice.

Selection remains partly non-neural and should remain so.  The reason is simple: self-evolution needs a grounding signal outside the model's preference for plausible text.  If the model both proposes and declares success without external verification, the loop can collapse into self-confirmation.  Selection therefore depends on unit tests, benchmarks, mathematical objective functions, execution traces, or separately controlled judges.  Even when an LLM judge is used, robust systems prefer low-temperature evaluation, multiple samples, calibration sets, or verifiable subcomponents.  The evaluator is the anchor that prevents an evolutionary language loop from becoming mere rhetoric.

AlphaEvolve is the most prominent large-scale example of the LLM-as-evolutionary-operator paradigm \cite{novikov2025alphaevolve}.  It is described as a Gemini-powered coding agent for scientific and algorithmic discovery.  Its architecture combines a fast model for breadth, Gemini Flash, with a stronger model for depth, Gemini Pro.  The practical intuition is that exploration and exploitation need different model economics.  Breadth requires many candidate ideas at acceptable cost and latency.  Depth requires more expensive reasoning for difficult improvements, abstraction, or repair.  A dual-model architecture therefore creates a search policy over model capabilities, not only over program candidates.

AlphaEvolve extends the FunSearch style from single evolved functions to broader code artifacts.  It uses automated evaluators and an evolutionary database.  The public descriptions emphasize quality-diversity search, especially MAP-Elites, rather than a single best-so-far chain.  Quality-diversity methods aim to illuminate a space of solutions by preserving high-quality artifacts across diverse behavioral niches.  In algorithm discovery, this matters because the first solution that reaches a local optimum may not be the stepping stone to the next breakthrough.  A slower or less elegant candidate may contain a structural idea that becomes valuable after later recombination.

MAP-Elites maintains an archive indexed by behavior descriptors.  Let $\mathcal{B}$ be a discretized descriptor space and let $b(a)\in\mathcal{B}$ map an artifact to a cell.  The archive stores at most one elite per cell:
\begin{equation}
  E[c] = \arg\max_{a: b(a)=c} F(a).
\end{equation}
When a child $a_c$ is evaluated, it is inserted if its cell is empty or if it improves the elite already in that cell:
\begin{equation}
  E[b(a_c)] \leftarrow
  \begin{cases}
    a_c, & E[b(a_c)] = \varnothing,\\
    a_c, & F(a_c) > F(E[b(a_c)]),\\
    E[b(a_c)], & \text{otherwise.}
  \end{cases}
\end{equation}
The archive is therefore not a leaderboard but a map of high-performing diversity.

\begin{algorithm}[t]
\caption{LLM-guided MAP-Elites for evolutionary code discovery}
\label{alg:llm-map-elites}
\begin{algorithmic}[1]
\REQUIRE Seed programs $P_0$, evaluator $\operatorname{Eval}$, descriptor function $b$, LLM ensemble $\mathcal{M}$, budget $T$
\STATE Initialize archive $E$ with empty cells
\FORALL{$p\in P_0$}
  \STATE $(F(p), b(p), r(p)) \leftarrow \operatorname{Eval}(p)$
  \STATE Insert $p$ into $E$ if it is elite for cell $b(p)$
\ENDFOR
\FOR{$t=1$ to $T$}
  \STATE Select one or more elites $P_t$ from $E$ using fitness, novelty, age, and coverage weights
  \STATE Construct an evolution prompt with parent code, scores, descriptors, failures, and instructions
  \STATE Choose model $m_t\in\mathcal{M}$, e.g., fast model for breadth or strong model for depth
  \STATE Sample candidate edits or complete child programs $C_t\sim m_t(\cdot)$
  \FORALL{$c\in C_t$}
    \STATE Parse, sanitize, and sandbox $c$
    \STATE $(F(c), b(c), r(c))\leftarrow \operatorname{Eval}(c)$
    \IF{$c$ is valid and improves the elite in cell $b(c)$}
      \STATE $E[b(c)]\leftarrow c$
    \ELSE
      \STATE Store failure trace $r(c)$ for diagnostics or future prompts
    \ENDIF
  \ENDFOR
\ENDFOR
\RETURN archive $E$ and the global best elite $\arg\max_{c\in E}F(c)$
\end{algorithmic}
\end{algorithm}

The algorithm box abstracts AlphaEvolve-like systems while also describing open implementations.  Several design choices are hidden inside the prompt constructor.  The prompt may include parents sorted by score, a concise failure analysis, a list of forbidden regressions, or a request for a patch rather than a full file.  It may ask the model to preserve function signatures, stay within an evolution region, or add only deterministic code.  It may also use model routing: a fast model proposes many rough edits, while a stronger model repairs, explains, or intensifies promising branches.  In this sense, the prompt is not merely an input string; it is an operator specification.

AlphaEvolve's reported achievements illustrate why this architecture attracted attention.  The local materials report that it found a 48-multiplication algorithm for $4\times4$ complex matrix multiplication, described as the first improvement over the Strassen-era baseline in 56 years.  They also report that across more than fifty mathematical problems, AlphaEvolve rediscovered state-of-the-art results for about 75\% and found better solutions for about 20\%.  In infrastructure applications, the same source reports improvements such as recovering 0.7\% of worldwide compute in data-center scheduling, simplifying hardware accelerator design, and speeding LLM training kernels by 23\% or 32\% in particular attention-related settings.  These numbers are domain-specific and should be read as examples of evaluator-grounded optimization rather than universal guarantees.

The dual-model architecture is a major systems lesson.  In evolutionary search, the marginal value of a model call depends on where it is used.  Early generations benefit from diversity and cheap exploration.  Later generations may require deeper reasoning to escape plateaus.  A system can allocate model budgets analogously to mutation rates or temperature schedules in classical evolution.  A cheap model can populate the archive with many variants; an expensive model can be invoked when a niche is promising, when a failure trace suggests a precise repair, or when a parent has high novelty.  Self-evolution therefore becomes a resource-allocation problem over artifacts, evaluators, and models.

Quality-diversity also reframes success.  A single best program is useful for deployment, but an archive of elites is useful for continued self-improvement.  If the environment changes, a non-best elite may become the best starting point.  If a later prompt needs examples of different trade-offs, the archive supplies them.  If human researchers want insight, the map reveals which behavioral niches are easy, hard, empty, or overrepresented.  This is particularly relevant for AI self-evolution because open-ended improvement depends on preserving stepping stones that do not immediately maximize the current score.

The main risks are evaluator overfitting and specification gaming.  If the evaluator measures speed on a narrow input distribution, evolution may discover brittle shortcuts.  If tests are visible or inferable, a model may produce code that passes tests without solving the intended task.  If performance is measured on a shared machine without careful isolation, noise may drive selection.  Systems therefore need hidden tests, randomized inputs, resource controls, multiple metrics, and audit trails.  They also need a policy for failed candidates: failures are valuable search information, but including too many raw failures in prompts can pollute context and encourage degenerate patterns.  A good archive stores both successes and summarized failure modes.

\subsection{Program Evolution and Mathematical Discovery: FunSearch}
\label{subsec:funsearch}

FunSearch is the canonical demonstration that LLM-guided program evolution can produce mathematical discoveries rather than only better benchmark prompts \cite{romeraparedes2023mathematical}.  The system combines a pretrained LLM with an evaluator and a population database.  Instead of searching directly over solutions, it searches over programs that generate or prioritize solutions.  This distinction is crucial.  A raw combinatorial object may be too large, brittle, or unstructured for efficient search.  A program can encode a reusable heuristic, a constructive rule, or a priority function that generalizes across instance sizes.  The LLM's knowledge of code and mathematics becomes useful because the candidate artifact is executable source text.

The FunSearch pipeline has three central components: a sampler, an evaluator, and a programs database.  The database produces prompts containing prior program variants.  The sampler asks the LLM to complete a new version of the function marked for evolution.  The evaluator parses and executes the candidate, computes a score, and registers successful programs back into the database.  Users mark the problem with two decorators: one function is the run/evaluation entry point and another function is the function to evolve.  The evolved function may be a priority function for greedily constructing a mathematical object or a heuristic used inside a solver.

The local implementation analysis reports that FunSearch uses a population structure with islands, clusters, and programs.  Islands are independent subpopulations.  Clusters group programs with the same score signature across test cases.  Within a cluster, shorter programs are preferred, applying an Occam-style bias toward concise heuristics.  Periodic island resets remove weaker islands and restart them from strong founders, which helps avoid premature convergence while preserving successful discoveries.  This is a classical evolutionary idea adapted to LLM-generated code.

The prompt design is one of FunSearch's most important contributions.  A prompt may show a chain of versions such as \texttt{priority\_v0}, \texttt{priority\_v1}, and an empty \texttt{priority\_v2} with a docstring saying it should improve the previous version.  Earlier versions appear with imports and helper functions, allowing the LLM to see the evolving code context.  Versions can be sorted so that stronger examples appear later, exploiting recency effects in transformer context.  The prompt does not merely request a solution; it stages an evolutionary lineage and invites the model to continue it.

FunSearch's evaluator is deliberately strict.  It checks that the generated code parses, runs within a sandbox and timeout, does not call ancestor functions directly, and returns a numeric score.  Candidate code may be trimmed with AST-based recovery when the LLM emits incomplete text, but semantic admission depends on execution.  The evaluator also records a signature: a vector of scores across test instances.  This signature is used for clustering, so two programs with the same aggregate score but different behavior can be distinguished if their instance-level patterns differ.

A simplified FunSearch loop can be described as:
\begin{algorithm}[t]
\caption{FunSearch-style program evolution}
\label{alg:funsearch-style}
\begin{algorithmic}[1]
\REQUIRE Initial program with \texttt{@run} evaluator and \texttt{@evolve} target; LLM; programs database $D$; evaluators $\mathcal{E}$
\STATE Initialize islands in $D$ with the seed program
\WHILE{budget remains}
  \STATE Sample an island and construct a prompt from high-scoring program versions
  \STATE Ask the LLM to complete the next evolved function version
  \STATE Parse candidate code; trim or reject malformed bodies
  \STATE Execute candidate in sandbox on benchmark instances
  \IF{candidate runs, avoids ancestor calls, and returns numeric scores}
    \STATE Compute signature vector and reduced score
    \STATE Register program in the corresponding island and cluster
  \ELSE
    \STATE Discard or log the failure
  \ENDIF
  \IF{reset period elapsed}
    \STATE Reinitialize weak islands from founders sampled among strong islands
  \ENDIF
\ENDWHILE
\RETURN best programs and their mathematical constructions
\end{algorithmic}
\end{algorithm}

FunSearch's headline achievements include new constructions for the cap set problem and better heuristics for bin packing.  It also discovered solutions in related combinatorial settings such as admissible sets and cyclic graph independent sets.  These results are scientifically significant because the output programs can be inspected.  A human can read the evolved priority function, reason about why it works, and potentially translate it into mathematical insight.  This inspectability differentiates program search from black-box neural prediction.

It is important to correct a common attribution confusion.  The first improvement over the Strassen-era result for $4\times4$ complex matrix multiplication is associated with AlphaEvolve in the local AlphaEvolve materials, not with the original FunSearch paper.  FunSearch is the predecessor that established LLM-plus-evaluator program evolution for mathematical discovery.  AlphaEvolve later extended the approach to larger code artifacts and reported the matrix-multiplication improvement.  For historical accuracy, Table~\ref{tab:funsearch-alphaevolve-results} lists FunSearch's discoveries separately from the AlphaEvolve successor result.

\begin{table*}[t]
\centering
\caption{Representative program-evolution discoveries.  The matrix-multiplication row is included because it is often discussed alongside FunSearch, but the reported result belongs to AlphaEvolve rather than the original FunSearch system.}
\label{tab:funsearch-alphaevolve-results}
{\TableTight
\begin{tabularx}{\textwidth}{@{}L{0.16\textwidth}L{0.19\textwidth}L{0.20\textwidth}Y@{}}
\toprule
\textbf{System} & \textbf{Domain} & \textbf{Reported result} & \textbf{Interpretation} \\
\midrule
FunSearch & Cap set constructions & Better construction than previously known in the studied setting & Demonstrates that evolved programs can yield new mathematical objects rather than merely solve training instances. \\
FunSearch & Bin packing & Improved heuristic rules & Shows that program evolution can discover practical algorithms with interpretable logic. \\
FunSearch & Cyclic graphs and related combinatorics & New constructive programs in selected problems & Supports the claim that function-space search generalizes across mathematical domains with evaluators. \\
AlphaEvolve & $4\times4$ complex matrix multiplication & 48 multiplications, reported as first improvement over Strassen-era baseline in 56 years & Successor system extends program evolution to broader algorithmic code and quality-diversity search. \\
AlphaEvolve & 50+ mathematical problems & Recovered SOTA for roughly 75\% and improved roughly 20\% in local materials & Indicates breadth when evaluator design is available, but results are problem- and budget-dependent. \\
\bottomrule
\end{tabularx}}
\end{table*}

The methodological lesson is that mathematical discovery becomes possible when three conditions align.  First, the domain must admit a compact executable representation.  Second, the evaluator must be cheaper than the value of the discovery and reliable enough to guide search.  Third, the search system must preserve diversity because mathematical breakthroughs often require stepping stones.  FunSearch satisfies these conditions for several combinatorial problems.  It does not solve domains where evaluation is ambiguous, expensive, or dependent on human taste.

FunSearch also reveals why program evolution is different from ordinary code generation.  A code-generation benchmark usually asks the model to solve a problem once.  Program evolution asks the model to participate in an iterative population process.  The model is not judged only by the first sample; it is judged by whether its samples can be selected, recombined through context, and improved over generations.  A mediocre sample may still be valuable if it introduces a structural motif that later variants exploit.  The unit of scientific productivity is therefore the search process, not an individual completion.

For AI self-evolution, FunSearch provides a template that can be lifted from functions to agents.  Replace a priority function with a planning policy, a tool-selection rule, or a memory-update function.  Replace the mathematical evaluator with a task benchmark.  Replace island clusters by agent behavior signatures.  The same loop becomes agent evolution.  The difficulty increases because agents have larger state spaces, more sources of nondeterminism, and more opportunities to exploit evaluator loopholes.  Nevertheless, the FunSearch architecture remains a foundational pattern: evolve executable artifacts, evaluate them automatically, preserve a population, and expose interpretable lineages.

\subsection{Open-source Implementations: OpenEvolve, FunSearch, CodeEvolve, and SE-Agent}
\label{subsec:open-source-implementations}

Open-source implementations matter because they reveal the engineering surface hidden behind research papers.  A paper may state that an evaluator executes programs, but a reusable system must decide how to isolate file systems, limit memory, parse model outputs, checkpoint populations, route among model providers, recover from failed subprocesses, and present results to users.  It must also choose a license and a reproducibility posture.  In self-evolution, these engineering details are not incidental; they define what kinds of artifacts can safely evolve.

FunSearch's released code is best viewed as a research framework.  It exposes the sampler, evaluator, program database, decorators, AST utilities, and island model, but it leaves the LLM interface and secure execution environment as components that users must implement or adapt.  This is appropriate for a Nature paper whose purpose is to demonstrate a scientific method.  The code makes the algorithm intelligible and auditable, but it is not a turnkey platform for arbitrary codebase evolution.

OpenEvolve presents itself as an open-source implementation inspired by AlphaEvolve.  The current repository describes it as an evolutionary coding agent that turns LLMs into autonomous code optimizers, supports Python library usage, works with OpenAI-compatible providers, and includes MAP-Elites quality-diversity, island-based populations, an LLM ensemble, and an artifact side channel for error feedback.  Its README reports examples such as circle packing, adaptive sorting, filter design, prompt evolution, and GPU-related optimization, and the repository lists an Apache-2.0 license.  This positions OpenEvolve as a practical bridge between research-grade program evolution and user-facing automated optimization.

CodeEvolve is another open-source evolutionary coding agent.  Its arXiv abstract describes an islands-based genetic algorithm with modular LLM orchestration, execution feedback, task-specific metrics, context-aware recombination, adaptive meta-prompting, and targeted refinement.  The repository README describes Apache-2.0 licensing, optional MAP-Elites integration, distributed islands, prompt populations, SEARCH/REPLACE diff application, resource-limited sandboxed evaluation, and support for open-weight and proprietary models through OpenAI-compatible APIs.  CodeEvolve is therefore especially useful as an engineering reference for multi-island orchestration, exploration--exploitation scheduling, and structured patch generation.

SE-Agent is different in substrate.  It evolves reasoning trajectories for multi-step LLM code agents rather than directly evolving only source functions.  Its repository describes three operations: Revision, Recombination, and Refinement.  Revision analyzes failed trajectories and generates alternative strategies.  Recombination synthesizes new trajectories from strengths of multiple paths.  Refinement removes redundancies and risk patterns from promising paths.  The repository reports 80\% Top-1 performance on SWE-bench Verified and lists an MIT license.  This makes SE-Agent a key example of evolutionary ideas applied to agent traces and problem-solving procedures, not only to code artifacts.

\begin{table*}[t]
\centering
\caption{Comparison of representative open and semi-open evolutionary code or trajectory systems.  Public repository details can change; the table records the evidence available in the local research materials and web verification performed for this draft.}
\label{tab:open-source-evolution-systems}
{\TableTight
\begin{tabularx}{\textwidth}{@{}L{0.12\textwidth}L{0.18\textwidth}L{0.34\textwidth}L{0.14\textwidth}Y@{}}
\toprule
\textbf{System} & \textbf{Evolved artifact} & \textbf{Search and evaluation mechanism} & \textbf{Status} & \textbf{Best use case} \\
\midrule
FunSearch & Single marked Python function inside a problem scaffold & Island model, program clusters by score signature, LLM completion of versioned functions, and sandboxed execution for cap set, bin packing, admissible/cyclic graph tasks. & Apache-2.0 code; CC-BY materials in local notes. & Interpretable mathematical program discovery. \\
OpenEvolve & Functions, programs, prompts, and optimization targets & MAP-Elites quality-diversity, islands, LLM ensemble, artifact feedback, and user-defined evaluators for circle packing, prompt optimization, sorting, and scientific computing. & Apache-2.0 repository. & Practical AlphaEvolve-style experimentation with provider flexibility. \\
CodeEvolve & Code regions and algorithmic solutions & Islands-based genetic algorithm, context-aware recombination, adaptive meta-prompting, optional MAP-Elites, and sandboxed task-specific metrics. & Apache-2.0 repository. & Reproducible multi-island code evolution and ablation studies. \\
SE-Agent & Multi-step reasoning trajectories for code agents & Revision, recombination, and refinement across trajectories, evaluated on SWE-bench Verified with repository-reported 80\% Top-1. & MIT repository. & Improving agent reasoning procedures over real software issues. \\
AlphaEvolve & Broader algorithmic code artifacts & Gemini Flash/Pro dual-model evolution with a quality-diversity archive for math problems, matrix multiplication, data-center scheduling, chip design, and LLM training kernels. & Not generally public in local notes. & Large-scale evaluator-rich algorithm discovery. \\
\bottomrule
\end{tabularx}}
\end{table*}

The comparison highlights three axes.  The first axis is artifact granularity.  FunSearch evolves one function; OpenEvolve and CodeEvolve can evolve larger code regions; SE-Agent evolves trajectories; AlphaEvolve targets broader algorithmic code.  Larger granularity increases expressive power but also increases failure modes.  A single function can be sandboxed and scored relatively cleanly.  A codebase-level agent may need dependency installation, compilation, integration tests, performance profiling, and security isolation.

The second axis is diversity management.  FunSearch uses islands and signature clusters.  OpenEvolve and CodeEvolve expose MAP-Elites or island-based populations.  SE-Agent uses trajectory pools and recombination.  These mechanisms all address the same problem: greedy best-only search loses stepping stones.  In self-evolving systems, diversity is not decorative.  It is the memory of alternative hypotheses about how to solve the task.

The third axis is evaluation.  FunSearch and AlphaEvolve are strongest when the evaluator is a crisp objective function.  OpenEvolve and CodeEvolve generalize by letting users define evaluators, which increases applicability but shifts responsibility to the practitioner.  SE-Agent's evaluator is a software-engineering benchmark where success depends on resolving real repository issues.  This is more realistic but also more expensive and harder to reproduce.  The more open-ended the artifact, the more evaluation becomes a systems problem.

Licensing also shapes adoption.  Apache-2.0 code in FunSearch, OpenEvolve, and CodeEvolve supports permissive use with patent grants, while SE-Agent's MIT license is similarly permissive but shorter.  AlphaEvolve, by contrast, is described in papers and reports but is not generally available as a full public system in the local notes.  This means that open-source successors are not merely clones; they are the artifacts through which the community tests which parts of the AlphaEvolve/FunSearch paradigm are reproducible outside a large industrial lab.

A practical takeaway is that open-source evolutionary systems need robust defaults.  Users should not have to design from scratch how to checkpoint an archive, how to summarize failure traces, how to separate training and validation evaluations, or how to prevent an evolved program from reading hidden answers.  The next generation of frameworks should treat these as first-class concerns.  Configuration should expose model routing, evaluator budgets, archive policy, diversity descriptors, sandbox limits, and validation gates.  Without these controls, self-evolution can appear to work in demos while failing under scientific replication.

\subsection{Agent Self-Evolution Systems}
\label{subsec:agent-self-evolution-systems}

The systems above evolve prompts, functions, code, or trajectories.  Agent self-evolution raises the abstraction level further: the evolving artifact is an agent architecture, a workflow, a tool interface, a memory system, or the agent's own source code.  This shift is conceptually important because agents are not single input--output functions.  They perform multi-step reasoning, call tools, maintain state, interact with environments, and may coordinate with other agents.  Evolution therefore has to optimize not only what the agent says, but how it senses, plans, acts, remembers, critiques, and modifies itself.

ADAS, or Automated Design of Agentic Systems, formalizes this direction \cite{hu2024adas}.  Its Meta Agent Search algorithm uses a meta agent to write code for new agent architectures.  The search space is the set of valid Python programs implementing agents, which is Turing-complete in principle.  Each iteration, the meta agent reviews previously discovered designs and their performance, writes a new agent, evaluates it on benchmark tasks, and records the result.  The output is the best design discovered in the archive.  This is evolutionary code discovery applied to agent design rather than to mathematical heuristics.

ADAS's significance lies in the scope of its search space.  Hand-designed agent frameworks usually expose a fixed menu: chain-of-thought, reflection, tool use, debate, tree search, or planning.  A code-level design search can combine these motifs or invent new control flows.  The local materials report that ADAS discovered agents outperforming state-of-the-art hand-designed agents on benchmarks such as ARC and Google Research Football, with transfer across domains and across model backbones.  The exact numbers are less detailed in the local summaries than for some other systems, but the qualitative result is that agent design itself becomes an optimization target.

The Darwin G{\"o}del Machine (DGM) extends this idea with open-ended archives and self-modification \cite{zhang2025darwingodel}.  It maintains a growing archive of agent implementations, samples parent agents with a diversity bias, uses a foundation model to modify the agent code, evaluates the child on coding benchmarks, and keeps improvements.  The local materials report improvement on SWE-bench from 20.0\% to 50.0\% and on a Polyglot multi-language coding benchmark from 14.2\% to 30.7\%.  DGM is ``Darwin'' because it uses population search and selection; it is ``G{\"o}del'' because agents can modify their own code in an open-ended process inspired by self-referential improvement.

DGM illustrates the value of an archive over a single lineage.  If a child fails, the parent remains.  If a branch discovers a useful tool-use pattern but not the best score, it can persist as a stepping stone.  If a later task distribution changes, the archive may contain a better starting point than the current champion.  This is more robust than hill-climbing and closer to open-ended evolution.  The trade-off is cost: each child must be evaluated on meaningful software tasks, and code modifications may break syntax, dependencies, or safety assumptions.

G{\"o}del Agent takes a more direct self-referential route \cite{yin2024godel}.  It allows an agent to inspect and modify its own runtime logic using LLM-driven monkey patching.  The loop is: execute on tasks, collect results and feedback, analyze the current code and behavior, propose a patch, apply it at runtime, validate, and accept or roll back.  This architecture is closer to online self-modification than to population-level evolution.  Its safety mechanisms include rollback, validation, and high-level objective constraints.  The local materials report improvements on HotPotQA, ALFWorld, and broader text tasks, though the summarized numbers are qualitative rather than a detailed table.

G{\"o}del Agent exposes a different risk profile.  Runtime monkey patching is flexible and fast, but it is Python-specific and can introduce nondeterminism.  It also blurs the boundary between a candidate under evaluation and the evaluator environment.  A population archive such as DGM can quarantine children; a runtime-patching agent must be especially careful about rollback and validation because the executing system is the object being changed.  For safety-critical self-evolution, this distinction matters.

Agent Symbolic Learning offers a more structured and less code-destructive view of agent self-evolution \cite{zhou2024symbolic}.  It treats an agent as a symbolic network whose learnable ``weights'' are prompts, tools, and pipeline connections.  After executing a task, an LLM produces a language loss from the trajectory.  Then language gradients are backpropagated through the symbolic network, and optimizers update prompt nodes, tool nodes, or pipeline nodes.  The formulas from the local materials are:
\begin{equation}
  \mathcal{L}_{\mathrm{lang}} = \operatorname{LLM}(\mathcal{P}_{\mathrm{loss}}(\tau)),
\end{equation}
\begin{equation}
  \nabla^n_{\mathrm{lang}} = \operatorname{LLM}\left(\mathcal{P}_{\mathrm{gradient}}(\nabla^{n+1}_{\mathrm{lang}}, \mathcal{I}_n, \mathcal{O}_n, \mathcal{P}_n, \mathcal{T}_n, \mathcal{L}_{\mathrm{lang}})\right),
\end{equation}
followed by symbolic updates to prompts, tools, and pipeline structure.  The metaphor is backpropagation in language space.  Unlike DGM, the update target is not arbitrary source code; it is a structured agent graph.

The benchmark numbers in the local materials suggest broad but moderate gains.  On HotPotQA, symbolic learning improves F1 from 36.2 to 39.1 and exact match from 22.8 to 25.6.  On MATH with a GPT-4 backbone, accuracy rises from 56.0\% to 60.7\%.  On HumanEval, pass@1 improves from 68.3\% to 73.2\%.  On a software-development executability score from 1 to 4, the reported average increases from 2.4 for agents to 3.8.  These results support the idea that language-space credit assignment can improve deployed agents without weight updates.

Absolute Zero moves self-evolution into reinforcement learning with verifiable rewards and no external data \cite{zhao2025absolutezero}.  A single model plays two roles: task proposer and task solver.  It generates reasoning tasks, attempts to solve them, and receives rewards from a code executor.  The proposer is rewarded for tasks that are solvable, challenging, and diverse; the solver is rewarded for correct solutions.  The local materials describe this as zero-data self-play RLVR, with PPO or GRPO-style updates for both roles.  The key equation is a dual reward loop:
\begin{equation}
  t \sim \pi_{\mathrm{prop}}(\cdot\mid s),\qquad y\sim\pi_{\mathrm{solve}}(\cdot\mid t),\qquad r=\operatorname{Execute}(y,t).
\end{equation}
The model improves by generating its own curriculum.

Absolute Zero is not evolutionary code search in the narrow sense, but it belongs in this chapter because it shares the self-evolution loop: generate tasks, evaluate outcomes, update the generator/solver, and repeat.  It also demonstrates the importance of verifiability.  Without a code executor, self-generated tasks could drift into ambiguous or self-serving problems.  With execution-based rewards, the system can train on an automatically grounded curriculum.  The local materials report that AZR-Coder-7B reaches state-of-the-art results among 7B models on coding averages and remains competitive on HumanEval+, MBPP+, LiveCodeBench, GSM8K, and MATH, despite using zero external training data.

SelfEvolve, from 2023, is an earlier code self-improvement pipeline \cite{jiang2023selfevolve}.  It has two stages.  First, the LLM generates knowledge such as API documentation or algorithm descriptions from its own parameters.  Second, generated code is executed in a sandbox, and error messages are fed back for iterative revision.  The local materials give formulas for knowledge-conditioned generation and revision, including:
\begin{equation}
  P(Y\mid K)=\prod_i p_\theta(Y_i\mid Y_{<i},X,K),
\end{equation}
where $K$ is self-generated knowledge, and a revision distribution conditioned on the original problem, generated code, knowledge, and error message.  SelfEvolve is closer to self-debugging than to population evolution, but it supplies a key mechanism later systems reuse: execution feedback as a natural-language training signal for the next attempt.

Together, these systems define a spectrum.  ADAS searches agent architectures in code space.  DGM evolves a population of self-modifying coding agents.  G{\"o}del Agent performs runtime self-modification.  Agent Symbolic Learning updates structured prompts, tools, and pipelines through language gradients.  Absolute Zero trains a model through self-generated verifiable tasks.  SelfEvolve iteratively refines code using self-generated knowledge and execution errors.  All instantiate self-evolution, but their update substrates differ: architecture code, agent source, runtime methods, symbolic graph weights, model policy, and generated code.

\begin{table*}[t]
\centering
\caption{Agent self-evolution systems by substrate, feedback signal, and reported benchmark evidence.}
\label{tab:agent-self-evolution-systems}
\begin{tabular}{p{0.14\linewidth}p{0.18\linewidth}p{0.20\linewidth}p{0.24\linewidth}p{0.14\linewidth}}
\hline
System & Evolved substrate & Feedback signal & Reported evidence in local materials & Main limitation \\
\hline
ADAS & Python programs implementing agent architectures & Benchmark scores on agent tasks & Outperforms hand-designed agents on ARC and GFootball; transfer across domains and models & Expensive search over huge code space \\
DGM & Archive of coding-agent implementations & SWE-bench and Polyglot scores & SWE-bench 20.0\% $\rightarrow$ 50.0\%; Polyglot 14.2\% $\rightarrow$ 30.7\% & Evaluation cost and safe self-modification \\
G{\"o}del Agent & Runtime agent logic via monkey patches & Task feedback plus validation & Reported gains on HotPotQA, ALFWorld, and multi-domain text tasks & Runtime patching risk and weaker portability \\
Agent Symbolic Learning & Prompts, tools, and pipeline connections & Language loss and language gradients & HotPotQA F1 36.2 $\rightarrow$ 39.1; MATH 56.0\% $\rightarrow$ 60.7\%; HumanEval 68.3\% $\rightarrow$ 73.2\% & Depends on strong LLM for gradient quality \\
Absolute Zero & Proposer/solver policy through RLVR & Code-executor rewards on self-generated tasks & 7B-level SOTA coding average in local notes; competitive math transfer & Curriculum collapse and executor-limited domains \\
SelfEvolve & Generated code plus self-generated knowledge & Sandbox errors and test feedback & DS-1000 49.4 $\rightarrow$ 57.2; HumanEval pass@1 74.39 $\rightarrow$ 85.98 & Mainly handles API/assertion errors; hand-written prompts \\
SE-Agent & Reasoning trajectories for code agents & SWE-bench Verified outcomes & Up to 55\% relative improvement in local notes; repository reports 80\% Top-1 & Expensive multi-trajectory search \\
\hline
\end{tabular}
\end{table*}

A unifying perspective is that agent self-evolution requires three memories.  The first is episodic memory: what happened in a particular task attempt.  The second is design memory: which prompts, tools, code changes, or trajectories have worked across attempts.  The third is evaluative memory: which measurements are trustworthy and under what conditions.  Systems that lack episodic memory repeat mistakes.  Systems that lack design memory cannot accumulate improvements.  Systems that lack evaluative memory overfit, regress, or reward hacks.

Another unifying issue is containment.  Prompt and symbolic updates are relatively easy to sandbox because they alter text fields.  Code updates need tests, permissions, dependency control, and rollback.  Runtime self-modification needs even stronger containment because the subject and mechanism of change overlap.  RL self-play needs reward containment because a model may learn to propose tasks that exploit the verifier.  Self-evolution is therefore as much an infrastructure problem as an algorithmic one.

\subsection{Benchmark Comparison and Evidence Quality}
\label{subsec:benchmark-comparison}

Benchmark comparison in this area is difficult because systems optimize different substrates under different budgets.  OPRO optimizes prompts and black-box variables.  FunSearch optimizes mathematical programs.  AlphaEvolve optimizes algorithms and infrastructure code.  DGM and SE-Agent target software-engineering agents.  Agent Symbolic Learning updates agent components across question answering, math, code, software development, and creative writing.  Absolute Zero trains models through self-play.  A single leaderboard would obscure these differences.  A better comparison asks: What is the artifact?  What is the evaluator?  What baseline is used?  How much search or training budget is consumed?  Is the result held out, reproducible, and resistant to evaluator gaming?

{\TableTight
\begin{longtable}{@{}L{0.13\textwidth}L{0.18\textwidth}L{0.17\textwidth}L{0.18\textwidth}L{0.28\textwidth}@{}}
\caption{Selected numeric results reported in the local research materials and verified web sources for open repositories where noted. Values should be compared within rows, not across unrelated benchmarks.}
\label{tab:numeric-benchmark-comparison}\\
\toprule
\textbf{System} & \textbf{Benchmark or domain} & \textbf{Baseline} & \textbf{Reported evolved result} & \textbf{Notes} \\
\midrule
\endfirsthead
\caption[]{Selected numeric benchmark results (continued).}\\
\toprule
\textbf{System} & \textbf{Benchmark or domain} & \textbf{Baseline} & \textbf{Reported evolved result} & \textbf{Notes} \\
\midrule
\endhead
\bottomrule
\endlastfoot
OPRO & Instruction optimization on reasoning tasks & Initial or prior prompts & Improved prompt scores over iterative search & Exact numbers depend on dataset and prompt placement; key evidence is the optimization trajectory. \\
OPRO & Linear regression/TSP optimization & Hand or heuristic initial candidates & Lower objective values through text-conditioned proposals & Demonstrates black-box optimization beyond prompts. \\
FunSearch & Cap set & Prior constructions in studied setting & Better evolved construction & Scientific discovery through executable program search. \\
FunSearch & Bin packing & Existing heuristic baselines & Improved heuristic program & Demonstrates practical heuristic discovery. \\
AlphaEvolve & $4\times4$ complex matrix multiplication & Strassen-era multiplication count & 48 multiplications & Reported as first improvement in 56 years in local materials. \\
AlphaEvolve & 50+ math problems & Known SOTA & Rediscovered SOTA for 75\%, improved 20\% & Indicates broad evaluator-driven search. \\
AlphaEvolve & Data-center scheduling & Existing Borg scheduling policy & 0.7\% worldwide compute recovered & Industrial result; evaluator and deployment details are domain-specific. \\
AlphaEvolve & LLM training kernels & Existing kernels & 23\% and 32\% speedups in reported attention settings & Performance depends on hardware and workload. \\
DGM & SWE-bench & 20.0\% & 50.0\% & Archive-based self-improving coding agents. \\
DGM & Polyglot coding & 14.2\% & 30.7\% & Multi-language transfer evidence. \\
Agent Symbolic Learning & HotPotQA & F1 36.2 / EM 22.8 & F1 39.1 / EM 25.6 & Structured language-gradient updates. \\
Agent Symbolic Learning & MATH & 56.0\% & 60.7\% & GPT-4 backbone in local materials. \\
Agent Symbolic Learning & HumanEval & 68.3\% pass@1 & 73.2\% pass@1 & Code-generation improvement without weight updates. \\
SelfEvolve & DS-1000 & ChatGPT 49.4 & 57.2 & Full pipeline, +15.8\% relative in local notes. \\
SelfEvolve & HumanEval & ChatGPT 74.39 pass@1 & 85.98 pass@1 & Approaches GPT-4 baseline in local table. \\
SelfEvolve & TransCoder & ChatGPT 88.3 accuracy / 88.5 pass@1 & 90.4 accuracy / 90.9 pass@1 & Translation and execution refinement. \\
Absolute Zero & 7B coding average & Data-trained comparison methods & SOTA among 7B models in local notes & Zero external data; exact aggregate depends on benchmark mix. \\
SE-Agent & SWE-bench Verified & Base code agents & Up to 55\% relative improvement; repository reports 80\% Top-1 & Trajectory-level evolution; web-verified repository claim. \\
OpenEvolve & Prompt optimization example & Baseline prompt & Repository reports +23\% HotpotQA accuracy in example & Open-source framework result, not a universal benchmark. \\
CodeEvolve & Algorithm discovery tasks & Open/closed baselines & arXiv abstract reports SOTA on several tasks & Requires paper-specific tables for exact per-task values. \\
\end{longtable}}

The table deliberately mixes precise and qualitative entries because the source materials vary in granularity.  For some systems, the local notes provide exact before/after values.  For others, they provide claims such as ``competitive,'' ``SOTA,'' or ``outperforms hand-designed agents'' without a row-by-row numeric breakdown.  A rigorous meta-analysis should return to the original papers, extract task definitions, budgets, seeds, and confidence intervals, and separate training, validation, and test performance.  The goal here is to map the evidence landscape rather than to declare a universal winner.

Several benchmark patterns recur.  First, the largest gains often occur when the baseline is a fixed agent or fixed prompt and the evolved system can exploit task-specific feedback.  This does not diminish the result; it clarifies that self-evolution is a form of adaptation.  Second, code and math tasks are overrepresented because they have verifiable rewards.  Third, software-engineering benchmarks such as SWE-bench are attractive because they test realistic agent behavior, but they are costly and sensitive to harness details.  Fourth, infrastructure optimization results can be highly valuable even when the percentage improvement looks small.  A 0.7\% compute recovery at global scale may be more economically significant than a large relative gain on a toy task.

Evidence quality depends on evaluator independence.  If candidate generation sees the exact tests, evolution may overfit.  If validation is performed on the same distribution as selection, performance may regress out of distribution.  If the evaluator is an LLM judge, model bias and prompt sensitivity enter the loop.  If the benchmark is public and widely used, contamination is possible.  Therefore, high-quality self-evolution evidence should report hidden tests, ablations, compute budgets, random seeds, failed-run rates, and regression checks.

Ablations are especially important.  In evolutionary systems, improvement can come from many sources: more model calls, better prompts, stronger base models, larger archives, smarter selection, more diverse parents, better evaluators, or simply more retries.  Without ablations, it is hard to know whether a proposed evolutionary mechanism is necessary.  For example, a MAP-Elites archive should be compared against best-only selection under the same evaluation budget.  A dual-model architecture should be compared against a single strong model with equal cost.  A trajectory recombination operator should be compared against independent retries and self-reflection.  The literature is beginning to include such ablations, but survey readers should treat headline scores as incomplete without them.

Another issue is budget normalization.  A system that evaluates ten thousand candidates may outperform a system that evaluates one hundred candidates for reasons unrelated to algorithmic design.  This is not a flaw if the use case permits the budget, but it must be reported.  Useful metrics include score versus number of evaluations, score versus wall-clock time, score versus token cost, and score versus human engineering time.  In industrial settings, the relevant metric may be return on compute investment.  In scientific discovery, the relevant metric may be novelty per evaluator call.

For agent self-evolution, transfer is a critical benchmark dimension.  If an evolved agent improves only on the tasks used for selection, it is a specialized optimizer.  If its design transfers to new tasks, models, or domains, it is closer to a general self-evolution method.  The local ADAS notes emphasize cross-domain and cross-model transfer.  DGM reports transfer across coding benchmarks.  Absolute Zero reports cross-domain transfer from code self-play to math reasoning.  Such results are important because self-evolution should ideally discover reusable mechanisms, not only benchmark-specific hacks.

\subsection{Design Patterns for Robust AI Self-Evolution}
\label{subsec:evolution-design-patterns}

Across the systems surveyed, several design patterns emerge.  The first is \textbf{representation discipline}.  The evolved artifact should have a clear boundary.  FunSearch marks one function.  CodeEvolve modifies code between markers with structured diffs.  Agent Symbolic Learning updates named prompts, tools, and pipeline nodes.  SE-Agent manipulates trajectories.  Clear boundaries make evaluation, rollback, and attribution possible.

The second pattern is \textbf{verifier-first design}.  The evaluator should be designed before large-scale generation begins.  A weak evaluator makes the system optimize the wrong target.  A slow evaluator makes search impractical.  A non-deterministic evaluator makes selection noisy.  A leaky evaluator invites reward hacking.  The most successful systems are those in which candidate quality can be measured automatically and cheaply enough to support many iterations.

The third pattern is \textbf{archive over memory}.  A chat transcript is not an archive.  An archive stores artifacts, scores, descriptors, provenance, failure traces, and selection status in a structured way.  It supports parent sampling, diversity analysis, rollback, and audit.  OPRO's history is a minimal archive.  FunSearch's programs database is a richer archive.  DGM's open-ended agent collection and AlphaEvolve's quality-diversity map are more powerful still.

The fourth pattern is \textbf{diversity preservation}.  Best-only loops are simple but brittle.  They discard variants that may become useful later.  Island models, MAP-Elites, signature clustering, and trajectory pools all preserve alternative hypotheses.  Diversity is not only about avoiding local optima; it is about preserving interpretability.  A diverse archive lets researchers compare families of solutions and identify mechanisms that correlate with success.

The fifth pattern is \textbf{failure utilization}.  In ordinary pipelines, failed generations are discarded.  In self-evolution, failures can be compressed into useful signals: syntax errors, timeout causes, failing test names, counterexamples, performance profiles, or reasoning blind spots.  SelfEvolve uses error messages for revision.  OpenEvolve-style artifact side channels feed diagnostics back to the LLM.  SE-Agent uses failed trajectories for revision.  The challenge is to summarize failures without flooding context.

The sixth pattern is \textbf{model-budget scheduling}.  Not every generation deserves the strongest model.  AlphaEvolve's Flash/Pro split suggests a general recipe: use cheaper models or higher temperatures for broad exploration, stronger models for deep repair or exploitation, and route calls based on archive state.  Open-source systems can implement analogous policies with open-weight and proprietary models.  Budget-aware evolution is essential for scaling beyond demonstrations.

The seventh pattern is \textbf{regression control}.  Any self-evolving system can forget or degrade.  Code changes may improve one metric while breaking another.  Prompt changes may overfit.  Agent patches may alter safety behavior.  Robust systems need baselines, validation gates, canary tasks, hidden tests, and rollback.  The acceptance criterion should be multi-objective when deployment requires reliability, not just maximum benchmark score.

The eighth pattern is \textbf{human-legible lineage}.  Scientific and engineering value increases when humans can inspect how a solution evolved.  FunSearch's versioned functions, DGM's archive, and CodeEvolve's genealogical tracking all support lineage analysis.  A lineage can reveal whether a result emerged from gradual refinement, a sudden imported idea, or evaluator exploitation.  This is important for trust and for converting evolved artifacts into human knowledge.

\subsection{Open Problems}
\label{subsec:evolution-open-problems}

The first open problem is evaluator generality.  Current successes concentrate in domains with automatic verification: code, math, games, scheduling, kernels, and benchmarked agents.  Many important AI tasks involve ambiguity, human preference, long-horizon social effects, or safety constraints that are hard to score.  Extending self-evolution to these domains requires richer evaluators, uncertainty estimates, and human-in-the-loop governance.

The second open problem is theoretical understanding.  LLM-driven mutation is not a simple random process.  It is conditioned on pretraining, prompt design, archive content, and model decoding.  Classical convergence analyses for evolutionary algorithms do not directly apply.  We need models of semantic variation, archive-induced bias, and evaluator noise.  Without theory, system design remains empirical and brittle.

The third open problem is contamination and benchmark exhaustion.  As systems repeatedly optimize public benchmarks, the benchmarks become less informative.  Self-evolution can consume a benchmark by discovering its quirks.  This is especially acute for software-engineering tasks where repositories, tests, and solutions may leak into model training data.  Dynamic benchmarks, private tests, and task generation with verifiable rewards are partial remedies.

The fourth open problem is safety of self-modifying code.  An evolved program may consume excessive resources, access unintended files, or learn to manipulate its evaluator.  Runtime self-modification raises additional concerns because the modification mechanism itself can be changed.  Practical systems need sandboxing, least-privilege execution, provenance tracking, policy checks, and emergency rollback.  Safety should be evaluated as a first-class metric rather than treated as an implementation detail.

The fifth open problem is multi-objective alignment.  Real deployments care about accuracy, cost, latency, maintainability, privacy, fairness, and robustness.  A scalar benchmark score cannot encode all of these.  MAP-Elites and Pareto archives offer one route, but descriptor choice is hard.  If the descriptors are poor, the archive may preserve diversity that is irrelevant while missing diversity that matters.

The sixth open problem is the relationship between self-evolution and weight learning.  OPRO, FunSearch, AlphaEvolve, DGM, and SE-Agent can improve frozen-model systems through external search.  Absolute Zero updates model policies through self-play RL.  Future systems may combine both: an external archive discovers artifacts, and a model is periodically trained or distilled on successful lineages.  The challenge is to retain verifiable gains without distilling evaluator hacks or reducing diversity.

The seventh open problem is scientific credit assignment.  When an evolved program solves a mathematical problem, which component deserves credit: the base model, the prompt, the evaluator, the archive, the human problem formulation, or the compute budget?  This is not only philosophical.  It affects reproducibility, intellectual property, and how researchers interpret discoveries.  Lineage-aware archives and detailed experiment logs are necessary for credible claims.

\subsection{Synthesis}
\label{subsec:evolution-synthesis}

Evolutionary code and algorithm discovery marks a transition in AI self-evolution from isolated generation to cumulative search.  OPRO shows that a natural-language history of candidates and scores can turn an LLM into a black-box optimizer.  FunSearch shows that the candidates can be executable programs whose evolved behavior yields mathematical discoveries.  AlphaEvolve shows that the same idea can scale with model ensembles, automated evaluators, and quality-diversity archives to broader algorithmic and infrastructure domains.  OpenEvolve and CodeEvolve show that the community is translating these ideas into reusable frameworks.  SE-Agent shows that evolutionary operations can act on reasoning trajectories.  ADAS, DGM, G{\"o}del Agent, Agent Symbolic Learning, Absolute Zero, and SelfEvolve show that agents themselves can become the evolving substrate.

The unifying equation is simple:
\begin{equation}
  \text{self-evolution} = \text{generation} + \text{evaluation} + \text{memory} + \text{selection} + \text{variation under constraints}.
\end{equation}
Each term is necessary.  Generation without evaluation is imagination.  Evaluation without memory is repeated trial and error.  Memory without selection is an unstructured log.  Selection without variation is stagnation.  Variation without constraints is unsafe.  The systems surveyed in this chapter differ in how they instantiate the terms, but they all rely on the same closed-loop principle.

For future survey and benchmark work, the most important recommendation is to report the entire evolutionary ecology.  Papers should specify the artifact representation, model routing, prompt construction, archive policy, diversity descriptors, evaluator details, validation split, compute budget, failure rate, and rollback policy.  Without these details, self-evolution results are difficult to compare and easy to overstate.  With them, the field can move from impressive demonstrations toward reliable, auditable, and reusable self-improving AI systems.


\subsection{System-by-System Analytical Notes}
\label{subsec:evolution-system-notes}

The following notes expand the comparison by applying a common analytical lens to each system.  They are included because evolutionary self-improvement methods are often described with similar vocabulary even when their engineering assumptions differ substantially.  For each system, the key questions are what evolves, what feedback is trusted, what memory is retained, and what failure mode dominates.

\paragraph{OPRO.} OPRO evolves instruction strings and serialized optimization variables.  Its trusted feedback signal is scored history, which means that the outer loop is only as reliable as the mechanism that assigns scores, errors, or rewards.  In the self-evolution taxonomy used in this chapter, OPRO should be read not as a standalone model capability but as an interaction pattern among a generator, an evaluator, and an archive.  The generator supplies candidate variation; the evaluator converts behavior into a selection signal; the archive determines which traces survive into future contexts.  This perspective is useful because it prevents a common misunderstanding: the LLM is not magically improving in isolation, but is participating in a larger cybernetic system that accumulates evidence across attempts.

For OPRO, the central engineering question is how much freedom the candidate artifact receives.  A narrow artifact boundary makes validation simpler and reduces the chance that the system changes its own measurement process.  A broad artifact boundary increases the chance of discovering qualitatively new strategies, but it also increases the burden on sandboxing, provenance, and regression tests.  The dominant risk is prompt overfitting and noisy scorer feedback.  A robust implementation should therefore include held-out validation, artifact lineage, failure summarization, and a rollback path.  In comparative studies, OPRO should be evaluated not only by its best score but also by its score as a function of evaluation budget, the diversity of retained candidates, and the stability of improvements under repeated seeds.

\paragraph{FunSearch.} FunSearch evolves versioned Python functions.  Its trusted feedback signal is sandboxed mathematical objective, which means that the outer loop is only as reliable as the mechanism that assigns scores, errors, or rewards.  In the self-evolution taxonomy used in this chapter, FunSearch should be read not as a standalone model capability but as an interaction pattern among a generator, an evaluator, and an archive.  The generator supplies candidate variation; the evaluator converts behavior into a selection signal; the archive determines which traces survive into future contexts.  This perspective is useful because it prevents a common misunderstanding: the LLM is not magically improving in isolation, but is participating in a larger cybernetic system that accumulates evidence across attempts.

For FunSearch, the central engineering question is how much freedom the candidate artifact receives.  A narrow artifact boundary makes validation simpler and reduces the chance that the system changes its own measurement process.  A broad artifact boundary increases the chance of discovering qualitatively new strategies, but it also increases the burden on sandboxing, provenance, and regression tests.  The dominant risk is evaluator leakage and premature convergence.  A robust implementation should therefore include held-out validation, artifact lineage, failure summarization, and a rollback path.  In comparative studies, FunSearch should be evaluated not only by its best score but also by its score as a function of evaluation budget, the diversity of retained candidates, and the stability of improvements under repeated seeds.

\paragraph{AlphaEvolve.} AlphaEvolve evolves algorithmic code artifacts.  Its trusted feedback signal is automated program and infrastructure evaluators, which means that the outer loop is only as reliable as the mechanism that assigns scores, errors, or rewards.  In the self-evolution taxonomy used in this chapter, AlphaEvolve should be read not as a standalone model capability but as an interaction pattern among a generator, an evaluator, and an archive.  The generator supplies candidate variation; the evaluator converts behavior into a selection signal; the archive determines which traces survive into future contexts.  This perspective is useful because it prevents a common misunderstanding: the LLM is not magically improving in isolation, but is participating in a larger cybernetic system that accumulates evidence across attempts.

For AlphaEvolve, the central engineering question is how much freedom the candidate artifact receives.  A narrow artifact boundary makes validation simpler and reduces the chance that the system changes its own measurement process.  A broad artifact boundary increases the chance of discovering qualitatively new strategies, but it also increases the burden on sandboxing, provenance, and regression tests.  The dominant risk is large compute budgets and limited public reproducibility.  A robust implementation should therefore include held-out validation, artifact lineage, failure summarization, and a rollback path.  In comparative studies, AlphaEvolve should be evaluated not only by its best score but also by its score as a function of evaluation budget, the diversity of retained candidates, and the stability of improvements under repeated seeds.

\paragraph{OpenEvolve.} OpenEvolve evolves user-defined code and prompts.  Its trusted feedback signal is custom evaluator callbacks and artifacts, which means that the outer loop is only as reliable as the mechanism that assigns scores, errors, or rewards.  In the self-evolution taxonomy used in this chapter, OpenEvolve should be read not as a standalone model capability but as an interaction pattern among a generator, an evaluator, and an archive.  The generator supplies candidate variation; the evaluator converts behavior into a selection signal; the archive determines which traces survive into future contexts.  This perspective is useful because it prevents a common misunderstanding: the LLM is not magically improving in isolation, but is participating in a larger cybernetic system that accumulates evidence across attempts.

For OpenEvolve, the central engineering question is how much freedom the candidate artifact receives.  A narrow artifact boundary makes validation simpler and reduces the chance that the system changes its own measurement process.  A broad artifact boundary increases the chance of discovering qualitatively new strategies, but it also increases the burden on sandboxing, provenance, and regression tests.  The dominant risk is quality depends on practitioner evaluator design.  A robust implementation should therefore include held-out validation, artifact lineage, failure summarization, and a rollback path.  In comparative studies, OpenEvolve should be evaluated not only by its best score but also by its score as a function of evaluation budget, the diversity of retained candidates, and the stability of improvements under repeated seeds.

\paragraph{CodeEvolve.} CodeEvolve evolves marked code regions and prompt populations.  Its trusted feedback signal is sandbox metrics with island orchestration, which means that the outer loop is only as reliable as the mechanism that assigns scores, errors, or rewards.  In the self-evolution taxonomy used in this chapter, CodeEvolve should be read not as a standalone model capability but as an interaction pattern among a generator, an evaluator, and an archive.  The generator supplies candidate variation; the evaluator converts behavior into a selection signal; the archive determines which traces survive into future contexts.  This perspective is useful because it prevents a common misunderstanding: the LLM is not magically improving in isolation, but is participating in a larger cybernetic system that accumulates evidence across attempts.

For CodeEvolve, the central engineering question is how much freedom the candidate artifact receives.  A narrow artifact boundary makes validation simpler and reduces the chance that the system changes its own measurement process.  A broad artifact boundary increases the chance of discovering qualitatively new strategies, but it also increases the burden on sandboxing, provenance, and regression tests.  The dominant risk is distributed execution complexity.  A robust implementation should therefore include held-out validation, artifact lineage, failure summarization, and a rollback path.  In comparative studies, CodeEvolve should be evaluated not only by its best score but also by its score as a function of evaluation budget, the diversity of retained candidates, and the stability of improvements under repeated seeds.

\paragraph{SE-Agent.} SE-Agent evolves multi-step reasoning trajectories.  Its trusted feedback signal is SWE-bench Verified issue resolution, which means that the outer loop is only as reliable as the mechanism that assigns scores, errors, or rewards.  In the self-evolution taxonomy used in this chapter, SE-Agent should be read not as a standalone model capability but as an interaction pattern among a generator, an evaluator, and an archive.  The generator supplies candidate variation; the evaluator converts behavior into a selection signal; the archive determines which traces survive into future contexts.  This perspective is useful because it prevents a common misunderstanding: the LLM is not magically improving in isolation, but is participating in a larger cybernetic system that accumulates evidence across attempts.

For SE-Agent, the central engineering question is how much freedom the candidate artifact receives.  A narrow artifact boundary makes validation simpler and reduces the chance that the system changes its own measurement process.  A broad artifact boundary increases the chance of discovering qualitatively new strategies, but it also increases the burden on sandboxing, provenance, and regression tests.  The dominant risk is cost of multi-trajectory exploration.  A robust implementation should therefore include held-out validation, artifact lineage, failure summarization, and a rollback path.  In comparative studies, SE-Agent should be evaluated not only by its best score but also by its score as a function of evaluation budget, the diversity of retained candidates, and the stability of improvements under repeated seeds.

\paragraph{ADAS.} ADAS evolves agent architectures written as Python programs.  Its trusted feedback signal is agent benchmark scores, which means that the outer loop is only as reliable as the mechanism that assigns scores, errors, or rewards.  In the self-evolution taxonomy used in this chapter, ADAS should be read not as a standalone model capability but as an interaction pattern among a generator, an evaluator, and an archive.  The generator supplies candidate variation; the evaluator converts behavior into a selection signal; the archive determines which traces survive into future contexts.  This perspective is useful because it prevents a common misunderstanding: the LLM is not magically improving in isolation, but is participating in a larger cybernetic system that accumulates evidence across attempts.

For ADAS, the central engineering question is how much freedom the candidate artifact receives.  A narrow artifact boundary makes validation simpler and reduces the chance that the system changes its own measurement process.  A broad artifact boundary increases the chance of discovering qualitatively new strategies, but it also increases the burden on sandboxing, provenance, and regression tests.  The dominant risk is vast Turing-complete search space.  A robust implementation should therefore include held-out validation, artifact lineage, failure summarization, and a rollback path.  In comparative studies, ADAS should be evaluated not only by its best score but also by its score as a function of evaluation budget, the diversity of retained candidates, and the stability of improvements under repeated seeds.

\paragraph{DGM.} DGM evolves archives of self-modifying coding agents.  Its trusted feedback signal is coding benchmark improvement, which means that the outer loop is only as reliable as the mechanism that assigns scores, errors, or rewards.  In the self-evolution taxonomy used in this chapter, DGM should be read not as a standalone model capability but as an interaction pattern among a generator, an evaluator, and an archive.  The generator supplies candidate variation; the evaluator converts behavior into a selection signal; the archive determines which traces survive into future contexts.  This perspective is useful because it prevents a common misunderstanding: the LLM is not magically improving in isolation, but is participating in a larger cybernetic system that accumulates evidence across attempts.

For DGM, the central engineering question is how much freedom the candidate artifact receives.  A narrow artifact boundary makes validation simpler and reduces the chance that the system changes its own measurement process.  A broad artifact boundary increases the chance of discovering qualitatively new strategies, but it also increases the burden on sandboxing, provenance, and regression tests.  The dominant risk is safe archive management and expensive evaluations.  A robust implementation should therefore include held-out validation, artifact lineage, failure summarization, and a rollback path.  In comparative studies, DGM should be evaluated not only by its best score but also by its score as a function of evaluation budget, the diversity of retained candidates, and the stability of improvements under repeated seeds.

\paragraph{Godel Agent.} Godel Agent evolves runtime methods and behavior logic.  Its trusted feedback signal is validation after monkey patching, which means that the outer loop is only as reliable as the mechanism that assigns scores, errors, or rewards.  In the self-evolution taxonomy used in this chapter, Godel Agent should be read not as a standalone model capability but as an interaction pattern among a generator, an evaluator, and an archive.  The generator supplies candidate variation; the evaluator converts behavior into a selection signal; the archive determines which traces survive into future contexts.  This perspective is useful because it prevents a common misunderstanding: the LLM is not magically improving in isolation, but is participating in a larger cybernetic system that accumulates evidence across attempts.

For Godel Agent, the central engineering question is how much freedom the candidate artifact receives.  A narrow artifact boundary makes validation simpler and reduces the chance that the system changes its own measurement process.  A broad artifact boundary increases the chance of discovering qualitatively new strategies, but it also increases the burden on sandboxing, provenance, and regression tests.  The dominant risk is runtime nondeterminism and rollback safety.  A robust implementation should therefore include held-out validation, artifact lineage, failure summarization, and a rollback path.  In comparative studies, Godel Agent should be evaluated not only by its best score but also by its score as a function of evaluation budget, the diversity of retained candidates, and the stability of improvements under repeated seeds.

\paragraph{Agent Symbolic Learning.} Agent Symbolic Learning evolves prompts, tools, and pipeline edges.  Its trusted feedback signal is language losses and gradients, which means that the outer loop is only as reliable as the mechanism that assigns scores, errors, or rewards.  In the self-evolution taxonomy used in this chapter, Agent Symbolic Learning should be read not as a standalone model capability but as an interaction pattern among a generator, an evaluator, and an archive.  The generator supplies candidate variation; the evaluator converts behavior into a selection signal; the archive determines which traces survive into future contexts.  This perspective is useful because it prevents a common misunderstanding: the LLM is not magically improving in isolation, but is participating in a larger cybernetic system that accumulates evidence across attempts.

For Agent Symbolic Learning, the central engineering question is how much freedom the candidate artifact receives.  A narrow artifact boundary makes validation simpler and reduces the chance that the system changes its own measurement process.  A broad artifact boundary increases the chance of discovering qualitatively new strategies, but it also increases the burden on sandboxing, provenance, and regression tests.  The dominant risk is dependence on strong LLM-generated gradients.  A robust implementation should therefore include held-out validation, artifact lineage, failure summarization, and a rollback path.  In comparative studies, Agent Symbolic Learning should be evaluated not only by its best score but also by its score as a function of evaluation budget, the diversity of retained candidates, and the stability of improvements under repeated seeds.

\paragraph{Absolute Zero.} Absolute Zero evolves self-generated tasks and solver policies.  Its trusted feedback signal is code-executor rewards, which means that the outer loop is only as reliable as the mechanism that assigns scores, errors, or rewards.  In the self-evolution taxonomy used in this chapter, Absolute Zero should be read not as a standalone model capability but as an interaction pattern among a generator, an evaluator, and an archive.  The generator supplies candidate variation; the evaluator converts behavior into a selection signal; the archive determines which traces survive into future contexts.  This perspective is useful because it prevents a common misunderstanding: the LLM is not magically improving in isolation, but is participating in a larger cybernetic system that accumulates evidence across attempts.

For Absolute Zero, the central engineering question is how much freedom the candidate artifact receives.  A narrow artifact boundary makes validation simpler and reduces the chance that the system changes its own measurement process.  A broad artifact boundary increases the chance of discovering qualitatively new strategies, but it also increases the burden on sandboxing, provenance, and regression tests.  The dominant risk is curriculum collapse and verifier scope.  A robust implementation should therefore include held-out validation, artifact lineage, failure summarization, and a rollback path.  In comparative studies, Absolute Zero should be evaluated not only by its best score but also by its score as a function of evaluation budget, the diversity of retained candidates, and the stability of improvements under repeated seeds.

\paragraph{SelfEvolve.} SelfEvolve evolves generated code with self-generated knowledge.  Its trusted feedback signal is sandbox errors and assertions, which means that the outer loop is only as reliable as the mechanism that assigns scores, errors, or rewards.  In the self-evolution taxonomy used in this chapter, SelfEvolve should be read not as a standalone model capability but as an interaction pattern among a generator, an evaluator, and an archive.  The generator supplies candidate variation; the evaluator converts behavior into a selection signal; the archive determines which traces survive into future contexts.  This perspective is useful because it prevents a common misunderstanding: the LLM is not magically improving in isolation, but is participating in a larger cybernetic system that accumulates evidence across attempts.

For SelfEvolve, the central engineering question is how much freedom the candidate artifact receives.  A narrow artifact boundary makes validation simpler and reduces the chance that the system changes its own measurement process.  A broad artifact boundary increases the chance of discovering qualitatively new strategies, but it also increases the burden on sandboxing, provenance, and regression tests.  The dominant risk is limited correction scope and hand-written prompts.  A robust implementation should therefore include held-out validation, artifact lineage, failure summarization, and a rollback path.  In comparative studies, SelfEvolve should be evaluated not only by its best score but also by its score as a function of evaluation budget, the diversity of retained candidates, and the stability of improvements under repeated seeds.


\subsection{Methodological Checklist for Reproducible Evolutionary Discovery}
\label{subsec:evolution-methodological-checklist}

A reproducible evolutionary discovery paper should describe the search process with the same care that a machine-learning paper describes data, model, and optimizer.  The first required item is the \emph{candidate schema}.  The paper should state whether the candidate is a prompt, a full program, a patch, a marked function, a trajectory, a tool graph, or a model-generated task.  It should also state which parts of the surrounding system are immutable.  Without this boundary, readers cannot distinguish an algorithmic improvement from an accidental change in the benchmark harness.  For code artifacts, the schema should include allowed imports, allowed file writes, time limits, memory limits, dependency versions, and whether the candidate can observe evaluator internals.

The second item is the \emph{proposal distribution}.  In LLM-guided evolution, the proposal distribution is not a simple Gaussian mutation.  It is a product of model choice, decoding temperature, prompt structure, parent sampling, context compression, and output parsing.  A reproducibility report should therefore include model names, model versions or dates, sampling parameters, number of samples per prompt, retry policy, prompt templates, and any post-processing rules.  If a system uses a dual-model policy, such as a cheap model for breadth and a stronger model for depth, it should report when each model is invoked and how model-routing decisions are made.  If failures are summarized back to the model, the summary template should be part of the method because it changes future variation.

The third item is the \emph{archive policy}.  A self-evolution system is defined by what it remembers.  Does it store every candidate, only valid candidates, only elites, or also failures?  Does it preserve lineages?  Are candidates clustered by behavior signatures, MAP-Elites cells, prompt embeddings, task domains, or hand-designed descriptors?  Does parent selection prefer recent candidates, high-fitness candidates, novel candidates, or underexplored niches?  These choices can dominate performance.  A best-only archive approximates hill climbing; an island archive approximates distributed evolutionary search; a MAP-Elites archive approximates landscape illumination.  Papers should report archive size, pruning rules, reset rules, and whether validation results can overwrite training scores.

The fourth item is the \emph{evaluator contract}.  The evaluator should be described as an interface with inputs, outputs, failure modes, and trust assumptions.  For mathematical program search, the evaluator may return a deterministic score.  For code agents, it may return pass/fail on a benchmark issue.  For agent trajectories, it may involve an execution environment and a judge.  The contract should define what counts as invalid, how timeouts are handled, whether partial credit is possible, and whether stochastic results are averaged.  Hidden tests and public tests should be separated.  If the evaluator is expensive, the paper should report how many evaluations were run and how many were discarded before scoring.

The fifth item is \emph{budget accounting}.  Evolutionary systems can look much stronger than single-shot baselines simply because they consume more attempts.  Fair comparison does not require identical budgets for all use cases, but it does require transparent budgets.  Authors should report token counts, model-call counts, candidate counts, valid-candidate counts, evaluator calls, wall-clock time, hardware, and monetary cost when possible.  The most informative plot is often performance versus evaluator budget rather than final score alone.  A system that reaches a moderate score quickly may be preferable to a system that reaches a higher score only after enormous search.

The sixth item is \emph{regression measurement}.  A candidate should not be accepted merely because it improves one scalar on one batch.  The system should measure whether it preserves previous capabilities, whether it generalizes to held-out tasks, and whether it remains safe under adversarial or randomized tests.  In code evolution, regression suites should include functional correctness, performance, determinism, and resource use.  In agent evolution, regression suites should include tool-use safety, refusal behavior where appropriate, memory consistency, and robustness to prompt perturbations.  The acceptance gate should be explicit.  For multi-objective systems, the paper should state whether it uses lexicographic ordering, scalarization, Pareto dominance, or niche-specific elites.

The seventh item is \emph{lineage reporting}.  A final artifact without a lineage is difficult to trust.  Lineage reporting should include the parent identifiers, mutation prompt, evaluation results, and reason for acceptance.  For scientific discoveries, it should include intermediate artifacts that reveal how the idea emerged.  This enables human researchers to inspect whether the system discovered a genuine mechanism or exploited an evaluator quirk.  Lineage is also useful for debugging: if a later generation regresses, the archive can identify the branch where the regression entered.

The eighth item is \emph{negative evidence}.  Evolutionary discovery papers often emphasize successful branches, but failed branches are scientifically valuable.  They reveal what the evaluator rejects, what the model tends to misunderstand, and which search operators are unproductive.  A concise failure taxonomy can be more useful than a long list of raw failed programs.  Typical categories include syntax errors, semantic type errors, timeout failures, test overfitting, numerical instability, nondeterminism, dependency misuse, unsafe file access, and reward hacking.  Reporting the distribution of failures helps other researchers reproduce the search and design better filters.

The ninth item is \emph{human intervention accounting}.  Some systems claim autonomous self-evolution, yet humans may adjust prompts, repair infrastructure, choose benchmarks, or manually inspect candidates during the run.  These interventions are not necessarily bad; they are often necessary for early research.  But they should be reported.  A clear distinction between autonomous search, human-curated search, and post-hoc human interpretation prevents overclaiming.  For deployment, the level of acceptable human intervention will depend on risk: scientific exploration may tolerate human curation, while autonomous production code modification requires stricter automation and oversight.

The tenth item is \emph{release artifacts}.  A reproducible release should include the final artifact, seed artifacts, prompt templates, evaluator code, benchmark splits, configuration files, and archive logs.  If a proprietary model prevents exact reproduction, the release can still provide enough structure for approximate reproduction with other models.  Open-source systems such as FunSearch, OpenEvolve, CodeEvolve, and SE-Agent are valuable because they expose this surrounding infrastructure, not merely because they publish a clever prompt.

\subsection{A Conceptual Taxonomy of Evolutionary Self-Improvement Loops}
\label{subsec:evolution-taxonomy-loops}

The systems in this chapter can be organized by the depth at which evolution intervenes.  At the shallowest level is \emph{instruction evolution}.  The artifact is a natural-language prompt, and the evaluator measures task performance.  OPRO is the representative case.  Instruction evolution is easy to deploy because it does not require code execution, but it is limited by the expressivity of the prompt and by the stability of LLM behavior under prompt changes.  It is best suited for tasks where a fixed model already has latent competence and needs better elicitation.

The next level is \emph{program-fragment evolution}.  The artifact is a function or marked code region.  FunSearch is the cleanest example.  This level has strong interpretability because the evolved fragment can be read and analyzed.  It also has strong containment because the evaluator can restrict the fragment's inputs and outputs.  The limitation is that the surrounding scaffold must already define the problem well.  If the key innovation requires changing the scaffold, a single-function search may be too narrow.

The third level is \emph{codebase evolution}.  The artifact may include multiple functions, files, or algorithmic modules.  AlphaEvolve, OpenEvolve, and CodeEvolve operate closer to this level.  Codebase evolution can discover larger design changes and optimize real systems, but it requires build systems, integration tests, performance harnesses, dependency management, and stronger sandboxing.  It also needs better credit assignment because a performance change may result from interactions among many edits.

The fourth level is \emph{trajectory evolution}.  The artifact is a sequence of thoughts, actions, tool calls, or intermediate plans.  SE-Agent is representative.  Trajectory evolution is attractive when the final code change is less important than the process by which an agent reaches it.  A revised or recombined trajectory can teach the agent how to approach similar problems.  The risk is that trajectory quality may depend heavily on the base model and environment state, making reproduction more difficult.

The fifth level is \emph{architecture evolution}.  The artifact is an agent design, workflow, or multi-agent organization.  ADAS searches this level by asking a meta agent to write agent programs.  DGM evolves agent implementations in an archive.  Architecture evolution can discover new reasoning and tool-use patterns, but it has the broadest search space and the hardest evaluation problem.  A good architecture may be robust across tasks, while a bad architecture may overfit to benchmark quirks in subtle ways.

The sixth level is \emph{self-referential runtime evolution}.  The artifact is the agent's own operating logic as it executes.  G{\"o}del Agent illustrates this level.  Runtime evolution is powerful because adaptation can happen online, but it requires strong rollback and containment.  The system must ensure that the mechanism for change does not corrupt the mechanism for evaluation.  This is the level at which safety concerns become most acute.

The seventh level is \emph{policy evolution through self-generated data}.  Absolute Zero represents this level.  The artifact being improved is no longer only external code or prompts; it is the model policy itself, trained through self-play tasks and verifiable rewards.  This level can produce more durable improvements because knowledge may enter the weights, but it also raises the cost and risk of training.  It requires reward design, curriculum control, and safeguards against self-generated distribution collapse.

These levels are not mutually exclusive.  A future self-evolving system may use OPRO-style prompt optimization to improve its mutation prompts, FunSearch-style program evolution to discover heuristics, MAP-Elites to preserve code diversity, SE-Agent-style trajectory recombination to improve problem-solving traces, ADAS-style architecture search to redesign the agent, and Absolute-Zero-style self-play to train the underlying model.  The research challenge is to compose these loops without creating uncontrolled feedback cycles.

\subsection{Implications for the Broader Survey}
\label{subsec:evolution-implications-survey}

For the broader topic of AI self-evolution, evolutionary code and algorithm discovery supplies the most concrete evidence that AI systems can accumulate improvements outside conventional supervised training.  The evidence is concrete because the artifacts are inspectable and the evaluators are often executable.  A prompt can be read, a function can be run, a patch can be tested, and an archive can be audited.  This makes the area a useful anchor for discussions that might otherwise become speculative.

The chapter also suggests that self-evolution should be studied as a systems discipline.  The LLM is important, but it is not the entire system.  The evaluator, archive, sandbox, prompt compiler, model router, benchmark harness, and logging infrastructure are equally important.  In many cases, the innovation lies in how these components are connected.  FunSearch's versioned prompt format, AlphaEvolve's quality-diversity archive, DGM's open-ended agent collection, and Agent Symbolic Learning's language-gradient graph are system-level contributions.

A second implication is that verifiability will shape the near-term geography of self-evolution.  Code, mathematics, theorem-like tasks, games, scheduling, chip design, and kernel optimization are likely to progress faster because they provide reliable feedback.  Domains such as negotiation, education, medicine, and social planning will require more careful evaluator design and governance.  The field should resist the temptation to generalize from verifiable domains to all forms of intelligence without qualification.

A third implication is that diversity is a safety and capability mechanism.  It is a capability mechanism because diverse archives preserve stepping stones and alternative strategies.  It is a safety mechanism because diversity makes it easier to compare artifacts and detect suspicious convergence toward evaluator hacks.  If every branch independently discovers the same shortcut, researchers should ask whether the shortcut reflects the intended objective.  If different branches solve the task through different mechanisms, humans gain more insight.

A fourth implication is that self-evolution changes the role of human researchers.  Humans may spend less time writing the final heuristic and more time designing the search space, evaluator, constraints, and audit process.  This is not a reduction in scientific responsibility.  It is a shift from direct construction to ecosystem design.  The quality of the resulting discovery depends on whether the ecosystem rewards genuine progress.

A fifth implication is that publication norms need to evolve.  A self-evolution result is not fully described by a final score and a method diagram.  It requires an archive summary, budget report, evaluator specification, lineage evidence, and failure analysis.  Without these, claims of autonomous improvement are hard to interpret.  With them, the community can compare systems on reliability, efficiency, and generality, not only on maximum benchmark score.

In summary, evolutionary code and algorithm discovery is the operational core of current AI self-evolution research.  It shows how LLMs can serve as optimizers, mutation operators, recombiners, debuggers, and meta-designers when embedded in a loop with memory and verification.  The field's next step is to make these loops more reproducible, safer, more budget-aware, and more general, while preserving the interpretability that makes program evolution scientifically valuable.
